% ═══════════════════════════════════════════════════════════════
% SPANISCH MUSTERLÖSUNG - Sequenz 07, Block b: Números 21-100
% ═══════════════════════════════════════════════════════════════
% Fachfarbe: #c41e3a (Karminrot)
% Lösungsfarbe: ROT
% ═══════════════════════════════════════════════════════════════

\documentclass[11pt, a4paper]{article}

% ═══════════════════════════════════════════════════════════════
% SW-MODUS TOGGLE
% ═══════════════════════════════════════════════════════════════
\newif\ifswmodus
\swmodusfalse    % Standard: Farbe

% ═══════════════════════════════════════════════════════════════
% PAKETE
% ═══════════════════════════════════════════════════════════════

\usepackage[utf8]{inputenc}
\usepackage[T1]{fontenc}
\usepackage[ngerman]{babel}
\usepackage{helvet}
\renewcommand{\familydefault}{\sfdefault}

\usepackage[margin=2cm]{geometry}
\usepackage{enumitem}
\usepackage{tabularx}
\usepackage{booktabs}
\usepackage{multirow}

\usepackage{xcolor}
\usepackage{framed}
\usepackage{tcolorbox}
\tcbuselibrary{skins}

\usepackage{fancyhdr}
\usepackage{lastpage}
\usepackage{needspace}

% ═══════════════════════════════════════════════════════════════
% FARBEN
% ═══════════════════════════════════════════════════════════════

\definecolor{spanisch-color}{HTML}{c41e3a}
\definecolor{grau-color}{HTML}{888888}
\definecolor{hellgrau-color}{HTML}{f5f5f5}
\definecolor{orange-color}{HTML}{b35c00}
\definecolor{loesung-color}{HTML}{c62828}

\definecolor{sw-dark}{HTML}{000000}
\definecolor{sw-gray}{HTML}{666666}
\definecolor{sw-white}{HTML}{ffffff}

\ifswmodus
    \colorlet{spanisch}{sw-dark}
    \colorlet{grau}{sw-gray}
    \colorlet{hellgrau}{sw-white}
    \colorlet{orange}{sw-dark}
    \colorlet{loesung}{sw-dark}
\else
    \colorlet{spanisch}{spanisch-color}
    \colorlet{grau}{grau-color}
    \colorlet{hellgrau}{hellgrau-color}
    \colorlet{orange}{orange-color}
    \colorlet{loesung}{loesung-color}
\fi

\newcommand{\fachfarbe}{spanisch}

% ═══════════════════════════════════════════════════════════════
% VARIABLEN
% ═══════════════════════════════════════════════════════════════

\newcommand{\thema}{Números 21-100 y precios}
\newcommand{\sequenz}{07}
\newcommand{\block}{b}
\newcommand{\fach}{Spanisch}
\newcommand{\klasse}{7}
\newcommand{\maxpunkte}{20}
\newcommand{\datum}{\today}

% ═══════════════════════════════════════════════════════════════
% KOPF- UND FUSSZEILE
% ═══════════════════════════════════════════════════════════════

\pagestyle{fancy}
\fancyhf{}
\renewcommand{\headrulewidth}{0.4pt}
\renewcommand{\footrulewidth}{0.4pt}

\lhead{\fach{} -- Klasse \klasse}
\chead{\textcolor{loesung}{\textbf{ML-\sequenz\block{} (MUSTERLÖSUNG)}}}
\rhead{\datum}

\lfoot{}
\cfoot{}
\rfoot{Seite \thepage\ von \pageref{LastPage}}

% ═══════════════════════════════════════════════════════════════
% BEFEHLE
% ═══════════════════════════════════════════════════════════════

\newcommand{\punktebox}[1]{\hfill\fbox{\small \_/#1}}

\newcommand{\loesung}[1]{\textcolor{loesung}{\textbf{#1}}}

\newcommand{\luecke}[1]{\rule{#1}{0.4pt}}

\newtcolorbox{merkkasten}[1][]{
    colback=hellgrau,
    colframe=\fachfarbe,
    colbacktitle=white,
    coltitle=\fachfarbe,
    boxrule=1pt,
    arc=0pt,
    left=8pt, right=8pt, top=8pt, bottom=8pt,
    fonttitle=\bfseries,
    attach boxed title to top left={yshift=-2mm, xshift=5mm},
    boxed title style={boxrule=0pt, colback=white},
    title=#1
}

\newtcolorbox{vokabelkasten}{
    colback=white,
    colframe=\fachfarbe,
    boxrule=0.5pt,
    arc=0pt,
    left=5pt, right=5pt, top=5pt, bottom=5pt
}

\newtcolorbox{hinweis}{
    colback=white,
    colframe=orange,
    colbacktitle=white,
    coltitle=orange,
    boxrule=1pt,
    arc=0pt,
    left=5pt, right=5pt, top=5pt, bottom=5pt,
    fonttitle=\bfseries,
    attach boxed title to top left={yshift=-2mm, xshift=5mm},
    boxed title style={boxrule=0pt, colback=white},
    title=Hinweis
}

\newenvironment{aufgabe}[1]{%
    \needspace{5\baselineskip}%
    \nopagebreak[4]%
    \vspace{10pt}
    \noindent\textbf{\textcolor{\fachfarbe}{Aufgabe #1}}
    \nopagebreak[4]%
    \vspace{5pt}
    \nopagebreak[4]%
    \begin{list}{}{\leftmargin=0pt}
    \item[]
}{%
    \end{list}
}

\newcommand{\es}[1]{\textcolor{\fachfarbe}{\textbf{#1}}}

% ═══════════════════════════════════════════════════════════════
% DOKUMENT
% ═══════════════════════════════════════════════════════════════

\begin{document}

% ═══════════════════════════════════════════════════════════════
% KOPFBEREICH
% ═══════════════════════════════════════════════════════════════

\begin{center}
    {\Large\bfseries\textcolor{\fachfarbe}{Arbeitsblatt: \thema}}\\[0.3em]
    {\large\textcolor{loesung}{\textbf{--- MUSTERLÖSUNG ---}}}
\end{center}

\vspace{5pt}

\noindent
\begin{tabularx}{\textwidth}{|X|X|X|}
\hline
\textbf{Name:} \textcolor{loesung}{Musterschüler/in} &
\textbf{Klasse:} \klasse &
\textbf{Datum:} \datum \\
\hline
\end{tabularx}

\vspace{15pt}

% ═══════════════════════════════════════════════════════════════
% MERKKASTEN 1: Zahlen 21-29
% ═══════════════════════════════════════════════════════════════

\begin{merkkasten}[Merkkasten 1: Die Zahlen 21-29]
\textbf{Achtung:} Die Zahlen 21-29 werden als \textbf{ein Wort} geschrieben!

\vspace{0.5em}

\begin{tabular}{llll}
21 = \loesung{veintiuno} & 22 = \loesung{veintidós} & 23 = \loesung{veintitrés} \\[0.8em]
24 = \loesung{veinticuatro} & 25 = \loesung{veinticinco} & 26 = \loesung{veintiséis} \\[0.8em]
27 = \loesung{veintisiete} & 28 = \loesung{veintiocho} & 29 = \loesung{veintinueve} \\
\end{tabular}

\end{merkkasten}

\vspace{10pt}

% ═══════════════════════════════════════════════════════════════
% MERKKASTEN 2: Zehner
% ═══════════════════════════════════════════════════════════════

\begin{merkkasten}[Merkkasten 2: Die Zehner 30-100]
\begin{tabular}{llll}
30 = \loesung{treinta} & 40 = \loesung{cuarenta} & 50 = \loesung{cincuenta} & 60 = \loesung{sesenta} \\[0.8em]
70 = \loesung{setenta} & 80 = \loesung{ochenta} & 90 = \loesung{noventa} & 100 = \loesung{cien} \\
\end{tabular}

\vspace{0.5em}

\textbf{Kombinationen:} Zehner + \es{y} + Einer

Beispiel: 35 = treinta \es{y} cinco
\end{merkkasten}

\vspace{10pt}

% ═══════════════════════════════════════════════════════════════
% AUFGABE 1: Zahlen zuordnen (mit Lösungen)
% ═══════════════════════════════════════════════════════════════

\begin{aufgabe}{1}
Verbinde die Zahl mit dem spanischen Wort. \punktebox{4}
\end{aufgabe}

\begin{center}
\begin{tabular}{rcl}
25 & \loesung{---} & sesenta y tres \loesung{(= 63)} \\[0.8em]
47 & \loesung{---} & veinticinco \loesung{(= 25)} \\[0.8em]
63 & \loesung{---} & ochenta y uno \loesung{(= 81)} \\[0.8em]
81 & \loesung{---} & cuarenta y siete \loesung{(= 47)} \\
\end{tabular}
\end{center}

\textit{\textcolor{grau}{Lösung: 25 $\leftrightarrow$ veinticinco, 47 $\leftrightarrow$ cuarenta y siete, 63 $\leftrightarrow$ sesenta y tres, 81 $\leftrightarrow$ ochenta y uno}}

\vspace{15pt}

% ═══════════════════════════════════════════════════════════════
% AUFGABE 2: Zahlen schreiben (mit Lösungen)
% ═══════════════════════════════════════════════════════════════

\begin{aufgabe}{2}
Schreibe die Zahlen auf Spanisch. \punktebox{4}
\end{aufgabe}

\noindent
a) 22 = \loesung{veintidós} \\[0.8em]
b) 38 = \loesung{treinta y ocho} \\[0.8em]
c) 54 = \loesung{cincuenta y cuatro} \\[0.8em]
d) 99 = \loesung{noventa y nueve} \\

\vspace{15pt}

% ═══════════════════════════════════════════════════════════════
% MERKKASTEN 3: Preise (mit Lösungen)
% ═══════════════════════════════════════════════════════════════

\begin{merkkasten}[Merkkasten 3: Los precios (Die Preise)]
\begin{tabular}{ll}
\es{¿Cuánto cuesta?} & = \loesung{Wieviel kostet es?} \\[0.8em]
\es{Cuesta X euros.} & = \loesung{Es kostet X Euro.} \\[0.8em]
\es{Cuesta X euros y Y céntimos.} & = \loesung{Es kostet X Euro und Y Cent.} \\
\end{tabular}
\end{merkkasten}

\vspace{10pt}

% ═══════════════════════════════════════════════════════════════
% AUFGABE 3: Preise lesen (mit Lösungen)
% ═══════════════════════════════════════════════════════════════

\begin{aufgabe}{3}
Lies die Preise und schreibe sie als Zahl. \punktebox{4}
\end{aufgabe}

\noindent
a) ``Cuesta veintitrés euros.'' = \loesung{23} € \\[0.8em]
b) ``Cuesta cuarenta y cinco euros.'' = \loesung{45} € \\[0.8em]
c) ``Cuesta setenta y ocho euros y cincuenta céntimos.'' = \loesung{78,50} € \\[0.8em]
d) ``Cuesta noventa y nueve euros.'' = \loesung{99} € \\

\vspace{15pt}

% ═══════════════════════════════════════════════════════════════
% AUFGABE 4: Preise sagen (mit Lösungen)
% ═══════════════════════════════════════════════════════════════

\begin{aufgabe}{4}
Schreibe die Preise auf Spanisch. \punktebox{4}
\end{aufgabe}

\noindent
a) 35 € = Cuesta \loesung{treinta y cinco} euros. \\[0.8em]
b) 67 € = Cuesta \loesung{sesenta y siete} euros. \\[0.8em]
c) 82,50 € = Cuesta \loesung{ochenta y dos} euros y \loesung{cincuenta} céntimos. \\[0.8em]
d) 100 € = Cuesta \loesung{cien} euros. \\

\vspace{15pt}

% ═══════════════════════════════════════════════════════════════
% AUFGABE 5: Dialog vervollständigen (mit Lösungen)
% ═══════════════════════════════════════════════════════════════

\begin{aufgabe}{5}
Vervollständige den Einkaufsdialog. \punktebox{4}
\end{aufgabe}

\begin{hinweis}
Verwende: ¿Cuánto cuesta? | Cuesta ... euros | Gracias | aquí tiene
\end{hinweis}

\vspace{0.5em}

\noindent
\textbf{Cliente:} ¡Buenos días! \loesung{¿Cuánto cuesta} la camiseta?

\textbf{Vendedor:} \loesung{Cuesta veinticinco euros.} \textit{(oder anderer Preis)}

\textbf{Cliente:} Vale. \loesung{Aquí tiene.}

\textbf{Vendedor:} Muy bien. ¡\loesung{Gracias}!

\vspace{0.5em}
\textit{\textcolor{grau}{Bewertung: ¿Cuánto cuesta? (1P), Preis (1P), Aquí tiene (1P), Gracias (1P)}}

\vspace{20pt}

% ═══════════════════════════════════════════════════════════════
% PUNKTEÜBERSICHT
% ═══════════════════════════════════════════════════════════════

\hrule
\vspace{10pt}

\begin{flushright}
\textbf{Punkte:} \loesung{\maxpunkte} / \maxpunkte
\end{flushright}

\end{document}
